\section{The sourced master identity and the interior charge}\label{sec:ward}

\subsection{A local current with the actual support}

In the sign convention of \cite[Lemma~4.8, Thm.~4.16]{RS}, the anomalous Master Ward identity for the one \(T\) gives
\begin{align}
 Q_0 S_T(V_f)
  &=\frac{i}{\hbar}\,T\bigl(e^{iV_f/\hbar},M_f\bigr),\nonumber\\
 M_f&=Q_0 V_f+\tfrac12\{V_f,V_f\}_{\rm std}
                  -i\hbar A_T(V_f).
 \label{eq:AWI}
\end{align}
Here braces denote the BV antibracket and \(A_T\) the local anomaly of this same extension of the time-ordered products. We use the coherent powers \(f^n\) because a constant-source germ \(f=a\) then equals the constant-coupling germ with coupling \(a\kappa\) at every homogeneous degree. This includes brackets and all local descendants with several vertices. Equal powers \(f\) on all distinct couplings would not have this property. Pure Einstein below means the full gauge-fixed BV Lagrangian, not only its graviton self-interaction; in particular the source jets of the antifield and ghost insertions are kept.

The classical part of the current is the negative of the first two terms of \(M_f\). This identifies its variational convention, not just its sign. In the generalized representative of \cite[Props.~4.3, 4.9]{RS}, the standard current differs from twice the source current by the contraction of the variational boundary potential (\(J_{\rm std}=2J_{\rm src}-\iota_Q\theta\), with that reference's test smearing). The nonminimal relative comparison also contains
\[
     B_0(V_f)=Q_0V_f-\{L_0,V_f\}_{\rm std}.
\]
Indeed its classically adjusted current
\(J_{\rm std}(L_0+V_f)-J_{\rm std}(L_0)-B_0(V_f)\)
is \(-Q_0V_f-\{V_f,V_f\}_{\rm std}/2\) by expansion of the bracket. We use this comparison classically; we do not apply its bare formula to quantum contact terms. The BV split, Euler contact terms and this boundary correction may not be dropped just because the profile is constant in the interior.

Here is an explicit descent of the full residual as a local functional of the switching profile. At a retained homogeneous coefficient of ghost number one, let \(m_f\) be a finite-jet local density for \(M_f\), including the coincident-vertex terms of \(A_T\). On the constant-jet subspace, \(m_f=f^n m_n\) at coupling degree \(n\). For any constant profile, the normalized unmarked QME holds as an integrated identity under compact off-shell field variations in an interior chart. Thus all Euler derivatives of \(m_n\) in the fluctuating BV fields are zero. Its field-independent part is zero by ghost number, with no external ghost constants. If \(\Theta_{m_n}\) is the boundary term in the first-variation formula for the density, scale the fluctuating fields \(\Phi\) by \(r\) and integrate:
\[
        m_n=\dd b_n,\qquad
        b_n(\Phi)=\int_0^1\Theta_{m_n}(r\Phi;\Phi)\,\dd r .
\]
A common covariant integration-by-parts convention uses the smooth background connection on the jet bundles. This does not assume that the density itself vanishes for constant \(f\).

Subtract \(\dd(f^n b_n)\) from \(m_f\). The difference vanishes when all positive-order source jets vanish. Write it as
\[
 P(f,\xi)-P(f,0)
     =\sum_{|I|\geq1}\xi_I\int_0^1
                      \partial_{\xi_I} P(f,r\xi)\,\dd r ,
       \qquad \xi_I=\nabla_I f,
\]
at the retained local finite-jet order. Integrate by parts on the chosen differentiated factor until one derivative of \(f\) remains; the other source jets in the coefficient stay in place. Thus in our convention
\begin{equation}
      J_f=-M_f=\int_M \dd f\wedge j[f] .
      \label{eq:descent}
\end{equation}
This is an off-shell identity of integrated local functionals; \(j[f]\) is a profile-dependent finite-jet three-form, not a universal profile-independent Noether current. No division by \(f\) or \(\dd f\) is involved. A marked derivative of \eqref{eq:descent} is taken in the same joint causal extension (the same \(T\) and Wick convention), not with each resulting graph separately finite-part renormalized. This is important for field-equation/propagator contacts. In particular \eqref{eq:descent} permits the genuine support split \(J_f=J_f^-+J_f^+\) of \cite[Remark~4.6, Lemma~5.1]{RS}.

For a buffered local insertion \(F\) supported inside the plateau, \cite[Eq.~(37)]{RS} then gives, before the dynamical quotient,
\begin{equation}
 Q_0 R_f(F)
   =R_f(\Db F)+\frac{i}{\hbar}[R_f(J_f^-),R_f(F)]_{\star},
 \qquad
 \Db F=Q_0F+\{V_f,F\}-i\hbar \Delta_{V_f}^{\rm ren}F .
 \label{eq:RS}
\end{equation}
The future term cancels by causal factorization; the past term reverses its order relative to time ordering, producing the commutator. Local coincident terms are retained in \(\Delta_{V_f}^{\rm ren}\). Equation~\eqref{eq:RS} for multilocal generators means the appropriate marked descendants of the same relative \(S\)-matrix, or the derivation on the algebra generated by the local on-shell insertions after the proof below. It is not an assumption that the ordinary pointwise product of interacting multilocal functionals obeys the same formula term by term.

\subsection{Noether insertion normalized inside the plateau}

Independently start with the gauge-fixed Einstein BRST Noether three-form \(j_{\rm N}\) from the variational boundary potential and divergence choice of \cite[App.~B]{BFR}. We use a compact bump one-form \(\eta_{\rm c}\) in a buffered collar with unit oriented integral along the time coordinate; compact spatial sections make this a compact current smear. Use a second, independent test \(\beta\) to mark the divergence within the plateau, not to turn off its interaction. The renormalized field-equation insertion and marked Ward relations in \cite[Eqs.~(40), (43), (67), (70)]{Rejzner} give the on-shell identity
\begin{equation}
    R_f\bigl((\dd j_{\rm N})(\beta)\bigr)
       = i\hbar R_f\bigl(A(V;\beta)\bigr),\qquad
    A(V;\beta)=\sum_{n}\Delta^{n}(V^{\otimes n};V(\beta)).
    \label{eq:markdiv}
\end{equation}
The notation in the marked slot means the infinitesimal local Lagrangian probe around the constant plateau; it is normalized there and is not a replacement of the coherent physical cutoff \(f\). All equalities can be understood on compact mark supports and their causal neighborhoods. We are \emph{not} imposing \(\Delta(V_f)=0\) on the switched presentation. The characteristic support of local anomalies means the coefficients in \eqref{eq:markdiv} depend only on the constant-coupling germ where the test or its derivatives are evaluated. Contributions outside the plateau belong to \eqref{eq:RS} instead.

To normalize the current insertion inside the plateau, first move derivatives off \(\beta\) in its local anomaly density. The anomaly is understood in the \emph{same nonminimal split} as \eqref{eq:AWI}, using the free-Ward-compatible time ordering (the tilded anomaly in \cite[Sec.~3.4, Eq.~(65)]{Rejzner}). A marked profile equal to one near the compactly supported variations of all BV fields reads the unmarked constant-coupling anomaly in that same family. The classical master equation holds separately for the fixed classical action, as an integrated germ identity. The constant-coupling QME normalized by \(T\) therefore implies that this anomaly density has vanishing Euler derivative \emph{off shell}. Source-gradient terms vanish only in this chosen interior germ. The field-scaling formula used above now gives a local three-form \(k_V\) with \(A(V;\beta)=(\dd k_V)(\beta)\) as integrated local functionals before the equation quotient. This is the jet version of eliminating the zeroth marked integral of a total-diagonal distribution; compare \cite[Sec.~4.6, Lemma~9]{Hollands}. The separate field-equation contacts in the actual current Ward identity are not discarded by this off-shell choice. Define \(j_{\rm Q}=j_{\rm N}-i\hbar k_V\). It is conserved in the common interacting on-shell algebra on the plateau, with the same prescription at every cut. It does not cancel the actual \(J_f\) in a region with profile gradients.

For completeness the charge implementation uses the \emph{differentiated} marked identity, not the false replacement \(\Delta_{V_f}^{\rm ren}F=0\). In the variation of \eqref{eq:markdiv} with a local test insertion \(F\) hold the primitive \(k_V\) at the physical \(V\) fixed in the independently specified current slot. The responses differentiating the retarded normalization of the undifferentiated \(A(V;\beta)\) and its subtraction \((\dd k_V)(\beta)\) cancel. The remaining local anomaly descendants are both
\begin{equation}
  \sum_n n\,\Delta^n(V^{\otimes(n-1)},F;V(\beta)),
   \qquad
  \sum_n\Delta^n(V^{\otimes n};F(\beta)),
  \label{eq:descendants}
\end{equation}
in addition to classical bracket and Euler terms. These are respectively the variation in an interaction slot and in the probe slot of the local anomaly. The retarded-minus-advanced (GLZ) identity takes two compact past/future mark profiles with difference \(1\) on a neighborhood of \(\supp F\) and derivative zero there. On the coincident distributions in \eqref{eq:descendants}, it replaces the marked terms precisely by the unmarked linearization \(\Delta_{V_f}^{\rm ren}F\), as in \cite[Eqs.~(43), (70)]{Rejzner}. All Euler terms are removed using the renormalized field-equation insertion identity after differentiation. This argument is local on the plateau and retains the two descendants, which is why a global zero-anomaly hypothesis on the switched \(V_f\) is unnecessary. The classical bracket contribution is fixed by the variational Noether identity; the resulting charge satisfies
\begin{equation}
 \Cq^{\rm N}=R_f\left(\int\eta_{\rm c}\wedge j_{\rm Q}\right),
 \qquad
 [R_f(\Db F)]_{\rm os}
  =\frac{i}{\hbar}[\Cq^{\rm N},R_f(F)]_{\star,{\rm os}} .
 \label{eq:inner}
\end{equation}
The conversion of the printed right bracket includes its degree: for an odd charge, \([Q,F]_\star=-(-1)^{|F|}[F,Q]_\star\). All operations here use the left convention fixed for \(Q_0\), not a degree-independent sign replacement. The commutator comparison with a switched defect parallels \cite[Thms.~15--17, Eq.~(102)]{TT}; we have specified its locality, variational reference and switched source here in a pure Einstein chart.

\subsection{On-shell equality, nilpotence, and smooth recutting}

On antifield-independent dynamical representatives the free equations remove \(\delta_0\) on the left of \eqref{eq:RS}; the remaining \(\gamma_0\) is the free charge commutator. Comparing \eqref{eq:RS} with \eqref{eq:inner} therefore gives
\begin{equation}\label{eq:center}
 [\Cq^{\rm N}-q_0+R_f(J_f^-),R_f(F)]_{\star,{\rm os}}=0 .
\end{equation}
Use the \emph{full gauge-fixed graded dynamical algebra} at this step, not only its gauge-invariant subalgebra. Algebraic multipliers are eliminated by their equations in the generating free Cauchy system. After that quotient its Green pairings are nondegenerate (indefinite is allowed). The standard smooth time-slice argument represents every smooth Cauchy datum by applying the free equation to a smoothly time-truncated solution near the slice. Thus a polynomial element central against all linear Cauchy generators is a scalar: its commutator with a linear field is exactly its Green directional derivative times \(i\hbar\), so all solution-direction derivatives vanish. This holds of graded Wick polynomials, including the anticommuting sectors. In \eqref{eq:center}, apply the argument to the first possible nonzero coupling coefficient, using the leading free term of each \(R_f(F)\) and dividing the free linear commutator by \(i\hbar\). Iterate to the retained order. There are no external odd constants in the dynamical algebra, and the difference has ghost number one; it vanishes. This proves \eqref{eq:keycharge} with \(\Cq=\Cq^{\rm N}\). The on-shell center cannot be replaced by an assertion of equality of off-shell currents: even the free current divergence contains the free Euler products.

The bulk BV differential is nilpotent on the local dynamical quotient by the normalized plateau QME. By the graded Jacobi identity \(\Cq\star\Cq\) commutes with this same generating algebra. Its ghost number two removes its possible central term by the same coefficient argument. In the real normalization the adjoint condition for the charge follows in the same way from that of its derivation (with the fixed Krein convention in the graded sector). All these operations are formal Wick operations on the retained degrees.

For two unit temporal smearings in the common plateau their difference is \(\dd\beta\) with compact primitive, so conservation of \(j_{\rm Q}\) gives equality of their charges. One can also present the same causal algebra using auxiliary cutoffs agreeing on an open comparison strip containing a full Cauchy slice. If the alteration \(W=W_-+W_+\) is respectively supported in the past and future of that strip, set \(U_-=S(V)^{-1}\star S(V+W_-)\). For a test \(F\) in the strip causal factorization gives
\[
 S_{V+W}(F)=U_-^{-1}\star S_V(F)\star U_- .
\]
The future piece cancels in the relative matrix, the past piece enters by conjugation. Differentiation in the current and observable slots identifies the same \(\Cq\), not \(q_0\) and \(R_f(J_f^-)\) separately. The field-equation time-slice property and these embeddings yield the same physical quotient maps with composition, or one may simply use a single \(V_f\) for every artificial cut. This is the statement of Theorem~\ref{thm:charge}; it supplies no equality between separately normalized boundary theories.

